\subsection{Regular-coherent coefficients and arbitrary valuation domains}
\label{C:regular-coherent-subsection}

We call a commutative ring \emph{regular coherent} if it is coherent
and every finitely presented module has a finite resolution by finite
projective modules.  We use the projective Grothendieck group of a
Schur algebra, rather than replacing the lattice exact category by
all finitely presented modules.

\begin{lemma}[Proved: generation through a perfect Schur lift]
\label{C:Schur-perfect-generation}
Let $k$ be a commutative ring, $n\geq1$, $B=k[\mathfrak S_n]$,
$V=k^n$, and
\[
 T=V^{\otimes n},\qquad S=\operatorname{End}_{B}(T),
\]
where $\mathfrak S_n$ acts on the right by place permutations.
Let $e\in S$ be the projection onto the multilinear weight summand
of $T$.  If $M$ is a left $B$-module finite projective over $k$,
put $F=Se\otimes_{eSe}M$, using $eSe\simeq B$.
If the underlying $k$-module of $F$ is perfect, then
$[M]\in R_k(\mathfrak S_n)$ is a $K_0(k)$-linear combination of the
Young permutation modules
\[
 k[\mathfrak S_n/\mathfrak S_\mu],\qquad
 \mathfrak S_\mu=\prod_j\mathfrak S_{\mu_j},\quad \mu\vdash n.
                                                               \tag{C.1.RC1}
\]
For every $k$ and $M$ as above, $F$ is finitely presented over $k$.
\end{lemma}
\begin{proof}
We first fix the structural inputs and their precise scope.
The algebra $S$ is the integral Schur algebra $S_k(n,n)$ and is
split quasi-hereditary over \emph{every} commutative ring
\cite[Theorem 4.8]{C-EfimovSchur}; the integral construction and
its standard filtrations are given in
\cite[Theorems 8.4.11 and 8.6.6]{C-KrauseSchur}.
Efimov's Proposition 3.16(2),(3), whose Section 3.1 explicitly
allows non-Noetherian rings, proves that $S$ is homologically
smooth over $k$ and that its standard modules form a full
exceptional collection in $\operatorname{Perf}(S)$
\cite{C-EfimovSchur}.
We spell out below the consequences needed for the lattice
exact category.  The Noetherian hypothesis in
\cite[Proposition 8.6.8]{C-KrauseSchur} is not discarded;
that proposition is not used.

\emph{The corner and its finite lift.}
The tensor-word basis of $T$ is permuted by $\mathfrak S_n$.
Its multilinear summand $T_\omega$ has the words using each
of the $n$ basis letters exactly once as basis, and is the
regular right $B$-module.  Its projection $e$ commutes with
place permutations.  Extension by zero from this summand gives
\[
 eSe=\operatorname{End}_B(T_\omega)
       \simeq\operatorname{End}_B(B_B)\simeq B,\qquad eF\simeq M.
                                                               \tag{C.1.RC2}
\]
The last identification uses left multiplication on $B_B$,
so it fixes the left and right module conventions.

The entries of a matrix in $S$ are constant on the orbits of
simultaneous place permutation on pairs of tensor words.
The indicator matrices of these finitely many orbits form a
$k$-basis, proving that $S$ is finite free and that its formation
commutes with arbitrary base change.
Consequently $Se$ and $eS$ are finite projective over $k$.
The module $F$ is the cokernel of
\[
 \begin{split}
 Se\otimes_k B\otimes_k M&\longrightarrow Se\otimes_k M,\\
 x\otimes b\otimes m&\longmapsto xb\otimes m-x\otimes bm.
 \end{split}                                                   \tag{C.1.RC3}
\]
Both displayed modules are finite projective over $k$.
Thus $F$ is finitely presented over $k$, without a coherence
assumption.

\emph{Perfectness and admissible resolutions.}
Homological smoothness says that the diagonal bimodule $S$ is
perfect over $S\otimes_k S^{\rm op}$; this is the ordinary tensor
product because $S$ is $k$-projective.
For a left $S$-complex $X$ perfect over $k$, apply
$-\otimes_S^{\mathbf L}X$ to a finite construction of this
bimodule from finite free bimodules, shifts, cones and direct
summands.  A free bimodule is sent to a finite sum of
$S\otimes_k^{\mathbf L}X$, which is perfect over $S$.
All the indicated operations preserve perfectness.
The diagonal is sent to $X$ itself, proving that $X$ is
$S$-perfect.  In particular the hypothesis on $F$ makes it
$S$-perfect.

A perfect module concentrated in degree zero has a finite
resolution by finite projective modules: starting with a
bounded projective complex representing it, split off its
acyclic positive-degree tail from the last term downwards
and augment the remaining complex to its zeroth cohomology.
Choose such a resolution
\[
 0\longrightarrow Q_r\longrightarrow\cdots\longrightarrow Q_0
       \longrightarrow F\longrightarrow0.
\]
The functor $e(-)$ is exact, so this gives a resolution of $M$.
Each $eQ_i$ is finite projective over $k$, since $Q_i$ is a
summand of a finite sum of copies of $S$ and $eS$ is
$k$-projective.  The surjection onto the $k$-projective module
$M$ splits over $k$, making its kernel finite projective.
Repeating this argument upwards proves that every short
exact sequence of this resolution is admissible in the
lattice category.  Hence
\[
 [M]=\sum_{i=0}^{r}(-1)^i[eQ_i]\quad\text{in }R_k(\mathfrak S_n).
                                                               \tag{C.1.RC4}
\]

\emph{Generation of the projective Grothendieck group.}
Write $\Delta_\lambda$ for the standard Schur modules.
Choose a linear refinement of dominance on the partitions.
The full exceptional collection in that order gives a semiorthogonal
decomposition with one component
$\Delta_\lambda\otimes_k\operatorname{Perf}(k)$ for each
partition $\lambda$ of $n$.
The projection coefficients are $k$-perfect: for a perfect
$S$-complex $X$, a finite projective resolution of
$\Delta_\lambda$ computes
$\mathbf R\operatorname{Hom}_S(\Delta_\lambda,X)$ by a
bounded construction from direct summands of finite sums of $X$.
Since $S$ is finite projective over $k$, $X$ is $k$-perfect,
and so are these coefficient complexes.
Its finite sequence of projection triangles therefore expresses every
class of $K_0(S)=K_0(\operatorname{Perf}(S))$ as a sum of
$[\Delta_\lambda]\,a_\lambda$, with $a_\lambda\in K_0(k)$.
This is a decomposition by actual perfect complexes and
projection triangles, so it includes direct summands and
does not infer generation merely from numerical characters.

For a partition $\mu$, let
\[
 P_\mu=\Gamma^{\mu_1}(V)\otimes_k\cdots
                   \otimes_k\Gamma^{\mu_\ell}(V).
\]
Here each polynomial functor is made an $S$-module by the
standard evaluation equivalence at $V=k^n$; see
\cite[Proposition 8.3.11]{C-KrauseSchur}.
In this equivalence the idempotent $e$ acts as the projection
to the multilinear weight space.
These are the divided-power projective Schur modules.
The integral Cauchy filtrations give
\[
 [P_\mu]=\sum_{\lambda\vdash n}K_{\lambda\mu}
             [\Delta_\lambda],
 \qquad K_{\lambda\mu}=0\ \text{unless }\lambda\unrhd\mu,
 \qquad K_{\mu\mu}=1.                                    \tag{C.1.RC5}
\]
Here the filtration terms and factors are finite free over
$k$; the proof over $\mathbf Z$ and its arbitrary base change
are included in the opened Theorem 8.4.11
\cite{C-KrauseSchur}.
Thus (C.1.RC5) holds over the present $k$.
The integer unitriangular matrix has an integer inverse.
It follows that the $[P_\mu]$ generate $K_0(S)$ as a
$K_0(k)$-module.

The functor from finite projective $S$-modules to the lattice
category given by $Q\mapsto eQ$ induces a
$K_0(k)$-linear map
\[
 K_0(S)\longrightarrow R_k(\mathfrak S_n).
\]
It sends the Euler class of $F$ to $[M]$ by (C.1.RC4).
For any standard module, its finite $S$-projective resolution
has $k$-projective last term $\Delta_\lambda$; the same
splitting argument shows that its image is
$[e\Delta_\lambda]$.
Consequently the generation just proved gives $[M]$ as a
$K_0(k)$-linear combination of $[eP_\mu]$.

We identify these last modules directly.
A divided-power monomial basis of $P_\mu$ in multilinear
weight records a distribution of the $n$ distinct basis
letters into ordered blocks of sizes $\mu_j$.
Permuting the basis letters permutes this basis transitively,
with stabilizer $\mathfrak S_\mu$.
On $T_\omega$ these permutations act by left regular
multiplication, so this is exactly the corner action in
(C.1.RC2).  Therefore
$eP_\mu\simeq k[\mathfrak S_n/\mathfrak S_\mu]$.
No factorial has been inverted, and no choice between
Specht and dual Specht conventions is required.
This proves all assertions.
\end{proof}

\begin{theorem}[Proved: the complete criterion for regular-coherent rings]
\label{C:regular-coherent-criterion}
Let $k\ne0$ be a commutative regular-coherent ring.
Then $\vartheta_k$ is an isomorphism if and only if $k$ has
characteristic zero and the unit map
$\mathbf Z\longrightarrow K_0(k)$ is an isomorphism.
No Noetherian or Krull-dimension hypothesis is required.
\end{theorem}
\begin{proof}
For any lattice $M$ in degree $n\geq1$, the module $F$ of
Lemma~\ref{C:Schur-perfect-generation} is finitely presented
over $k$ and hence perfect by regular coherence.
That lemma expresses $[M]$ as a $K_0(k)$-linear combination
of Young permutation modules.
If $K_0(k)=\mathbf Z$ by the unit map, these coefficients
are integers times $[k]$.
The permutation module for a partition $\mu$ is exactly
\[
 [k[\mathfrak S_n/\mathfrak S_\mu]]
       =\rho_{\mu_1}\cdots\rho_{\mu_\ell}
\]
under the induction product.  Thus $\vartheta_k$ is
surjective in every degree, with the degree-zero assertion
being the hypothesis on $K_0(k)$.
This argument does not use a noetherian support filtration
or identify the lattice group with $G_0(k[\mathfrak S_n])$.

For characteristic-zero $k$, the algebra
$k\otimes_{\mathbf Z}\mathbf Q$ is nonzero: otherwise some
nonzero integer would annihilate $1\in k$.
Choose a characteristic-zero residue field $L$ of this algebra.
Base change to $L$ preserves every admissible sequence
because the underlying $k$-sequence splits, and it sends
each $\rho_i$ to the corresponding generator over $L$.
The composite of $\vartheta_k$ with this base change is
the field isomorphism of Proposition~\ref{C:Qalgebras}.
Therefore $\vartheta_k$ is injective.

Conversely, its degree-zero map is
$\mathbf Z\to K_0(k)$, which must be an isomorphism.
Positive characteristic is excluded by
Proposition~\ref{C:positive_characteristic}, whose statement
includes arbitrary nonzero commutative rings of composite
positive characteristic.  These are precisely the two
claimed necessary conditions.
\end{proof}

\begin{corollary}[Proved: all valuation domains]
\label{C:valuation-domain-criterion}
Let $V$ be a valuation domain, with no restriction on its
rank or value group.  Then $\vartheta_V$ is an isomorphism
if and only if $V$ has characteristic zero.
In particular the conclusion holds for every nondiscrete
mixed-characteristic valuation domain.
\end{corollary}
\begin{proof}
The ring $V$ is local, so every finite projective module
is finite free and the rank identifies $K_0(V)$ with
$\mathbf Z$ via the unit map.
It is coherent: every nonzero finitely generated ideal
is principal because principal ideals are linearly
ordered, and is free because $V$ is a domain.

Every finite presentation matrix over $V$ can be
diagonalized by invertible row and column operations.
Among its finitely many nonzero entries choose one
dividing all the others, move it to the upper-left
corner, and use its multiples to clear its row and
column.  Repeat on the remaining smaller matrix.
Consequently every finitely presented $V$-module is
a finite direct sum of free modules and modules
$V/(a)$ with $a\ne0$.
The latter have the finite free resolutions
\[
 0\longrightarrow V\xrightarrow{\ a\ }V
        \longrightarrow V/(a)\longrightarrow0.
\]
Hence $V$ is regular coherent.  Apply
Theorem~\ref{C:regular-coherent-criterion}.
Neither the matrix argument nor its conclusion
requires a discrete or finite-rank valuation.
\end{proof}


\begin{theorem}[Proved: coherent regularity only at modular primes]
\label{C:coherent-regular-modular-locus}
Let $k\ne0$ be a coherent commutative ring of characteristic zero
such that $\mathbf Z\xrightarrow{\sim}K_0(k)$, and fix $n\geq1$.
If $k_{\mathfrak p}$ is regular coherent for every prime
$\mathfrak p$ containing $n!$, then the degree-$n$ map
$\vartheta_k$ is an isomorphism.
Consequently the entire map is an isomorphism if this local
regular-coherence condition holds at every point of positive
residue characteristic.  No regularity hypothesis is imposed
on points of residue characteristic zero.
\end{theorem}
\begin{proof}
Take a $k$-projective lattice $M$ and its finitely presented
Schur lift $F$ from Lemma~\ref{C:Schur-perfect-generation}.
Put $t=n!$.
After inverting $t$, the quotient map
\[
 Se\otimes_k M\longrightarrow Se\otimes_B M=F
\]
has the well-defined $k_t$-linear section
\[
 x\otimes_B m\longmapsto
       \frac1t\sum_{g\in\mathfrak S_n}xg^{-1}\otimes gm.
                                                               \tag{C.1.RC6}
\]
To check balancedness, replacing $x$ by $xh$ and
reindexing $g=uh$ changes the sum to the expression for
$x\otimes_B hm$.  Composing with the quotient returns
$x\otimes_B m$.  Thus $F_t$ is a direct summand of a finite
projective $k_t$-module and is finite projective.
At a prime containing $t$, regular coherence of
$k_{\mathfrak p}$ makes $F_{\mathfrak p}$ perfect.
Hence $F$ has finite projective dimension at every prime.

We verify that a uniform finite projective resolution follows
in the present coherent setting.
Choose a resolution of $F$ by finite free $k$-modules.
Coherence and finite presentation of $F$ ensure inductively
that all its syzygies are finitely presented; we set $Z_0=F$
and use $0\to Z_{r+1}\to k^{a_r}\to Z_r\to0$.
Let
\[
 U_r=\{\mathfrak p:(Z_r)_{\mathfrak p}
                       \text{ is projective over }k_{\mathfrak p}\}.
\]
These sets are open.  Indeed a finitely presented module
which is free at a prime has a finite basis there; the
basis and its inverse maps extend after inverting one
element outside that prime.  Their finitely many
identities can be made valid after one further such
localization, giving a free module on a neighborhood.
Projectivity over the local ring is freeness.
The exact syzygy sequences show $U_r\subseteq U_{r+1}$,
since a surjection onto a projective module splits.
Local finite projective dimension implies that the $U_r$
cover $\operatorname{Spec}k$.
Quasi-compactness supplies a single $r$ with
$U_r=\operatorname{Spec}k$.
A finitely presented module which is locally free at every
prime is finite projective.  Thus $Z_r$ is finite projective,
and the truncated resolution proves that $F$ is perfect.

Lemma~\ref{C:Schur-perfect-generation} now gives generation
by Young permutation modules with $K_0(k)$ coefficients.
The unit hypothesis makes these integer coefficients,
proving degree-$n$ surjectivity.  The characteristic-zero
residue-field argument in
Theorem~\ref{C:regular-coherent-criterion} proves injectivity.
For the final assertion, every prime containing $n!$ has
positive residue characteristic for $n\geq2$; the case
$n=1$, and degree zero, are exactly $K_0(k)=\mathbf Z$.
This proves the entire statement.
\end{proof}
